
\documentclass[11pt]{article}
\usepackage[margin=0.8in]{geometry}
\usepackage{amsmath,amssymb,booktabs,array,longtable}
\usepackage{graphicx}
\usepackage{hyperref}
\usepackage{enumitem}
\usepackage{titlesec}
\usepackage{float}

\usepackage{caption}
\usepackage{longtable}
\usepackage{pdflscape}
\setlist{nosep}

\begin{document}

\begin{center}
{\Large\textbf{Shiv Nadar University Chennai}}\\
\textbf{CS3807 -- Deep Learning Laboratory}\\
\textbf{Experiment 3}\\
{\large\textbf{Implementation of Convolutional Neural Networks (CNNs) for Image Classification}}
\end{center}

\renewcommand{\arraystretch}{1.25}

\begin{tabular}{|p{3.8cm}|p{8cm}|p{2cm}|}
\hline
Degree \& Branch & B.Tech Artificial Intelligence \& Data Science & Semester V\\ \hline
Subject Code \& Name & CS3807 -- Deep Learning Laboratory & AY: 2026--27\\ \hline
\end{tabular}

\section*{1. Objective}
To design, implement, and train a Convolutional Neural Network (CNN) using TensorFlow/Keras for multi-class image classification on the CIFAR-10 dataset. To analyze feature extraction, evaluate model performance using standard metrics, and visualize learned feature maps and filters to understand the CNN's learning process.

\section*{2. Description}
This experiment aims to provide hands-on experience in designing and implementing Convolutional Neural Networks (CNNs) for image classification tasks using the TensorFlow/Keras deep learning framework. The experiment begins with understanding the fundamental building blocks of CNNs, including convolutional layers, activation functions, pooling layers, flattening, and fully connected layers. Students investigate how convolution kernels extract meaningful visual features such as edges, textures, and patterns from images, while pooling operations reduce the spatial dimensions of feature maps and improve computational efficiency. The experiment also explores the influence of important CNN hyperparameters such as kernel size, stride, padding, number of filters, optimizer, batch size, and training epochs on model performance. Students visualize intermediate feature maps produced by convolutional layers to understand how CNNs progressively learn hierarchical image representations. Furthermore, they compare the effects of different pooling strategies and analyze the trainable parameters of convolutional layers. Using the CIFAR-10 image dataset containing ten object categories, students construct, train, and evaluate a CNN architecture for multi-class image classification. The performance of the model is assessed using evaluation metrics including accuracy, precision, recall, F1-score, confusion matrix, and classification report. Training and validation accuracy/loss curves, class distribution, sample images, and feature map visualizations are generated to interpret the learning behavior of the network. Through this experiment, students gain practical knowledge of CNN architecture design, feature extraction, model training, performance evaluation, and the application of deep learning techniques to computer vision problems.



\section*{3. Dataset}

\textbf{Dataset:} CIFAR-10

Training Images: 50,000 \\
Testing Images: 10,000 \\
Classes: 10 \\
Image Size: $32 \times 32$ pixels \\
Color Channels: 3 (RGB)



\section*{4. Experimental Procedure}

\textbf{Task 1: Dataset Loading and Exploration}

\begin{itemize}
    \item Load the CIFAR-10 image classification dataset using TensorFlow/Keras.
    \item Print the dataset dimensions and class labels.
    \item Display sample images from each class.
    \item Plot the class distribution of the dataset.
\end{itemize}


\bigskip

\textbf{Task 2: Data Preprocessing}

\begin{itemize}
    \item Normalize pixel values to the range $[0,1]$.
    \item Convert class labels to categorical format (one-hot encoding).
    \item Split the dataset into training and testing sets.
    \item Print the shapes of the processed datasets.
\end{itemize}


\bigskip

\textbf{Task 3: CNN Model Construction}

Construct a Convolutional Neural Network consisting of:

\begin{itemize}
    \item Convolutional layers with ReLU activation.
    \item Max Pooling layers for spatial dimensionality reduction.
    \item Flatten layer.
    \item Fully Connected (Dense) layer.
    \item Output layer with Softmax activation for 10-class classification.
\end{itemize}

Display the network architecture using \texttt{model.summary()}.


\bigskip

\textbf{Task 4: CNN Training}

\begin{itemize}
    \item Compile the CNN using an appropriate optimizer and categorical cross-entropy loss.
    \item Train the model for a specified number of epochs using mini-batch gradient descent.
    \item Record training and validation accuracy and loss after each epoch.
\end{itemize}

\bigskip

\textbf{Task 5: Hyperparameter Analysis}

\begin{itemize}
    \item Experiment with different kernel sizes, number of filters, padding, and stride values.
    \item Compare different pooling strategies such as Max Pooling and Average Pooling.
    \item Analyze the effect of optimizer, batch size, and number of training epochs on model performance.
\end{itemize}


\bigskip

\textbf{Task 6: Feature Map Visualization}

\begin{itemize}
    \item Extract feature maps from intermediate convolutional layers.
    \item Visualize the learned feature representations.
    \item Observe how low-level and high-level image features are progressively extracted.
\end{itemize}


\bigskip

\textbf{Task 7: Model Evaluation}

\begin{itemize}
    \item Evaluate the trained CNN on the test dataset.
    \item Compute Accuracy, Precision, Recall, and F1-Score.
    \item Generate the Confusion Matrix and Classification Report.
\end{itemize}


\bigskip

\textbf{Task 8: Performance Visualization}

\begin{itemize}
    \item Plot training and validation accuracy curves.
    \item Plot training and validation loss curves.
    \item Display correctly and incorrectly classified sample images.
    \item Summarize the overall CNN performance for the CIFAR-10 dataset.
\end{itemize}


\section*{5. Source Code}


The complete source code, datasets, generated plots, and report files for this experiment are available in the following GitHub repository:



\noindent\textbf{GitHub Repository:} \\
\url{https://github.com/VIKASNI-S/Deep-Learning-Lab/tree/main/Lab-3}\\


\section*{6. Plots and their Inference}



\begin{figure}[H]
\centering
\includegraphics[width=0.9\linewidth]{EPS_Figures/sample_images.eps}
\caption{Sample images from the CIFAR-10 dataset representing the ten object classes.}
\label{fig:sample_images}
\end{figure}
Sample images confirm the diversity of objects present in dataset.
\begin{figure}[H]
\centering
\includegraphics[width=0.8\linewidth]{EPS_Figures/class_distribution.eps}
\caption{Distribution of images across the ten classes in the CIFAR-10 training dataset.}
\label{fig:class_distribution}
\end{figure}
Dataset contains equal no.of images of all 10 classes --> balanced dataset. This prevents the model from being biased towads a particular class.


\begin{figure}[H]
\centering
\includegraphics[width=0.8\linewidth]{EPS_Figures/accuracy_curve.eps}
\caption{Training and validation accuracy over training epochs.}
\label{fig:accuracy}
\end{figure}
Training and Validation accuracy increase steadily with epochs, demonstrating that the CNN learns meaningful image features. The small gap depicts good generalization with minimal overfitting.

\begin{figure}[H]
\centering
\includegraphics[width=0.8\linewidth]{EPS_Figures/loss_curve.eps}
\caption{Training and validation loss over training epochs.}
\label{fig:loss}
\end{figure}
The loss decreases consistently during training, showing successful optimization of the network parameters. The validation loss follows a similar trend, indicating stable learning and good convergence.

\begin{figure}[H]
\centering
\includegraphics[width=\linewidth]{EPS_Figures/ConvLayer_1_FeatureMaps.eps}
\caption{Feature maps generated by the first convolutional layer.}
\label{fig:conv1}
\end{figure}
1st Convolutional layer capatues low level features such as edges and corners. They serve as basis for the deeper feature learning.

\begin{figure}[H]
\centering
\includegraphics[width=\linewidth]{EPS_Figures/ConvLayer_2_FeatureMaps.eps}
\caption{Feature maps generated by the second convolutional layer.}
\label{fig:conv2}
\end{figure}
2nd Convolutional laye extracts more complex featues from the previous layer. It identifies object parts and distinctive features.

\begin{figure}[H]
    \centering
    \includegraphics[
        width=0.9\linewidth,
        height=0.75\textheight,
        keepaspectratio
    ]{EPS_Figures/ConvLayer_3_FeatureMaps.eps}
    \caption{Feature maps generated by the third convolutional layer.}
    \label{fig:conv3}
\end{figure}
3rd Convolution layer learns high level semantic features like object shapes and class specific patterns. Also suppresses irrelevant background information.

\begin{figure}[H]
\centering
\includegraphics[width=\linewidth]{EPS_Figures/Learned_Filters_ConvLayer_1.eps}
\caption{Learned convolution kernels of the first convolutional layer.}
\label{fig:filters1}
\end{figure}
The learned kernels exhibit simple edge- and gradient-detecting patterns, indicating that the network has learned fundamental visual filters. These kernels form the basis for extracting low-level image features from the input.

\begin{figure}[H]
\centering
\includegraphics[width=\linewidth]{EPS_Figures/Learned_Filters_ConvLayer_2.eps}
\caption{Learned convolution kernels of the second convolutional layer.}
\label{fig:filters2}
\end{figure}
The kernels become more structured and complex than those in the first layer, reflecting the learning of texture- and pattern-oriented filters.


\begin{figure}[H]
    \centering
    \includegraphics[
        width=0.9\linewidth,
        height=0.75\textheight,
        keepaspectratio
    ]{EPS_Figures/Learned_Filters_ConvLayer_3.eps}
    \caption{Learned convolution kernels of the third convolutional layer.}
    \label{fig:filters3}
\end{figure}
The learned kernels appear increasingly abstract and less visually interpretable, as they are optimized to capture high-level semantic information. These filters contribute to distinguishing different object classes during classification.


\begin{figure}[H]
\centering
\includegraphics[width=0.85\linewidth]{EPS_Figures/confusion_matrix.eps}
\caption{Confusion matrix of the CNN model on the CIFAR-10 test dataset.}
\label{fig:confusion}
\end{figure}
Most predictions lie along the diagonal showing that the model correctly classifies the majority of images. There is chance for misclassification between visually similar classes.

\begin{figure}[H]
\centering
\includegraphics[width=0.8\linewidth]{EPS_Figures/hyperparameter_accuracy.eps}
\caption{Validation accuracy obtained using different optimizers and batch sizes.}
\label{fig:hyperparameter}
\end{figure}
The model accuracy varies with different hyperparameter settings, demonstrating that hyperparameter selection significantly influences CNN performance. The optimal combination of learning rate, batch size, optimizer, and number of epochs achieved the highest classification accuracy.


\section*{7. Results}

\subsection{Evaluation Metrics}

\begin{table}[H]
\centering
\caption{Performance Evaluation Metrics of the CNN Model}
\label{tab:evaluation_metrics}
\begin{tabular}{|l|c|}
\hline
\textbf{Metric} & \textbf{Value} \\ \hline
Accuracy  &  0.7472 \\ \hline
Precision &   0.7508\\ \hline
Recall    &  0.7472 \\ \hline
F1-Score  & 0.7461\\ \hline
\end{tabular}
\end{table}


\section*{8. Classification Report}

\begin{table}[H]
\centering
\caption{Classification Report of the CNN Model on the CIFAR-10 Test Dataset}
\label{tab:classification_report}
\resizebox{\textwidth}{!}{
\begin{tabular}{|l|c|c|c|c|}
\hline
\textbf{Class} & \textbf{Precision} & \textbf{Recall} & \textbf{F1-Score} & \textbf{Support} \\ \hline
Airplane   & 0.82 & 0.78 & 0.80 & 1000 \\ \hline
Automobile & 0.90 & 0.79 & 0.84 & 1000 \\ \hline
Bird       & 0.76 & 0.57 & 0.65 & 1000 \\ \hline
Cat        & 0.56 & 0.59 & 0.57 & 1000 \\ \hline
Deer       & 0.69 & 0.71 & 0.70 & 1000 \\ \hline
Dog        & 0.66 & 0.64 & 0.65 & 1000 \\ \hline
Frog       & 0.79 & 0.82 & 0.80 & 1000 \\ \hline
Horse      & 0.73 & 0.84 & 0.78 & 1000 \\ \hline
Ship       & 0.86 & 0.86 & 0.86 & 1000 \\ \hline
Truck      & 0.74 & 0.88 & 0.80 & 1000 \\ \hline
\textbf{Accuracy} & \multicolumn{3}{c|}{0.75} & 10000 \\ \hline
\textbf{Macro Average} & 0.75 & 0.75 & 0.75 & 10000 \\ \hline
\textbf{Weighted Average} & 0.75 & 0.75 & 0.75 & 10000 \\ \hline
\end{tabular}
}
\end{table}

\section*{9. References}
\begin{enumerate}
\item Goodfellow et al., \emph{Deep Learning}.

\item Alex Krizhevsky, ``The CIFAR-10 Dataset,'' Canadian Institute for Advanced Research (CIFAR), 2009.
\item TensorFlow/Keras Documentation.
\end{enumerate}


\end{document}
